\section{Export and Runtime}
\label{sec:deployment}
\subsection{Exported Graph and Strict Compilation}
The exported ONNX graph uses opset 20 and a static $1\times3\times640\times640$
input. Inspection finds 89 convolutions, 76 sigmoid and 76 multiplication nodes,
9 splits, 14 additions, 20 concatenations, 3 max-pooling operations, 2 resizes,
and one softmax. There are no MatMul, Gemm, Reshape, Transpose, Div, or Slice
operators. The sigmoid/multiplication pairs represent retained SiLU activations;
operator counts before fusion are not execution-kernel counts.

Compilation targets DLA core 0 in FP16 and calibrated INT8 with GPU fallback
disabled and explicit native I/O formats. Successful compilation is checked against both the verbose
placement report and the optimized-engine inspector. For each strict optimized engine,
the inspector reports a single DLA node corresponding to one DLA loadable; this
does not mean that the source graph contains one operation. Build and profile
records identify engine name, compiler version, and output bindings.

This placement verification is configuration-specific. TensorRT's DLA softmax support depends on
the target and shape, and the compiler warns that the selected approximation
can introduce numerical error \cite{nvidia_dla_working}. A source-level
rewrite that compiles on this Orin/TensorRT combination is not sufficient
evidence for another DLA generation, input size, or precision. Build success
also does not validate raw-output decoding.

\subsection{Native Layouts and the Host Boundary}
\label{sec:packing}
The FP16 strict build requests \texttt{fp16:dla\_hwc4} input and
\texttt{fp16:chw16} outputs. The input pads three channels to four in
interleaved storage, occupying $4\times640\times640\times2=3{,}276{,}800$
bytes. Output channels (144) are divisible by 16, requiring no channel padding.
The INT8 strict engine instead uses \texttt{int8:dla\_hwc4} input and
\texttt{int8:chw32} outputs. Its input occupies 1,638,400 bytes; CHW32 pads
144 channels to 160, totaling 1,344,000 output bytes across the three levels.
Both input layouts satisfy the documented row alignment
\cite{nvidia_dla_working}. GPU and stock DLA+GPU fallback engines retain linear FP16
external bindings in both precision modes.

For native INT8 bindings, the host quantizes normalized pixels as
$q=\operatorname{clip}(\operatorname{round}(x/s),-128,127)$ and dequantizes
raw outputs as $y=sq$. Per-tensor scales are read from the model's calibration
cache and checked against the explicit build-time I/O ranges; engine-hash-bound
metadata supplies the runtime scales. Thus, changing precision changes both
arithmetic and the host I/O path, not just the serialized engine.

The reference runner performs packing, unpacking, and INT8 conversion on the
CPU. Timed input stages include packing and host-to-device (H2D) transfer;
output stages include device-to-host (D2H) transfer, unpacking, and any
dequantization. These are not pure layout-conversion timings. A CUDA stream
and device-visible buffers remain necessary without GPU-assigned learned layers.
Logical shape alone is insufficient, so formats, vectorized dimensions,
component counts, and context strides are recorded. For the FP16 engine,
TensorRT reports CHW4 input and CHW16 outputs; the padded input interpretation
is validated for this shape, not for arbitrary CHW4 bindings.

\subsection{Confidence-First Host Decoding}
\label{sec:decode}
The host splits each $P_\ell$ into 64 regression channels and 80 class channels.
For confidence threshold $\tau\in(0,1)$, an anchor survives if
\begin{equation}
 \sigma\!\left(\max_c S_{\ell,c}(h,w)\right) > \tau.
 \label{eq:threshold}
\end{equation}
The implementation applies sigmoid before the comparison to preserve dense
floating-point threshold behavior. Regression decoding is independent of
class confidence and is postponed until after rejection. For each
survivor and box side $q\in\{l,t,r,b\}$, the DFL expectation is
\begin{equation}
 \hat d_q=\sum_{j=0}^{r-1}j\,
 \frac{\exp R_{qj}}{\sum_{k=0}^{r-1}\exp R_{qk}}.
 \label{eq:dfl}
\end{equation}
The decoder places cell centers at $(w+0.5,h+0.5)$, forms the four box
coordinates from these distances, scales by the level stride, applies class
sigmoid, and performs class-aware NMS. The final coordinates are mapped back
through the image's letterbox transform.

If $m$ of $N=8400$ anchors survive, dense class scanning remains
$O(Nn_c)$, while the regression softmax/expectation work falls from $O(4rN)$
to $O(4rm)$. This does not reduce raw-output transfer volume. In the FP16 strict
10\,Hz test, surviving anchors range from 0 to 254, with median 28.5 over the
100-image cycle. Identical threshold strictness, class selection, NMS settings,
and tie handling are necessary for dense/sparse agreement.
